% GENERATED FILE. Edit canonical lessons, then run npm run generate:books.
% canonical-source-hash: fnv1a64:867dd70481a50ae5
% canonical-lessons: ZH-C130-bingren, ZH-W130-bing-ill, ZH-C130-ganmao, ZH-C130-fashao, ZH-C130-teng, ZH-C130-kanbing

\chapter{Feeling Ill}
\label{ch:zh-feeling-ill}
\hlchaptermodality{130}

\begin{chapteropening}
\textbf{By the end of this chapter:} \emph{I can say I am ill, have a cold or a fever, or that something hurts, and go to the doctor.}

Complete the last lesson of chapter 130: I can say I am ill, have a cold or a fever, or that something hurts, and go to the doctor.
\end{chapteropening}

\section[\texorpdfstring{\zh{病人} (bìngrén)}{病人 (bìngrén)}]{\zh{病人} (bìngrén) — patient; sick person}

\label{lesson:ZH-C130-bingren}



\begin{quote}
Before the new one: say the Mandarin for convenient, then the Mandarin for interesting.
\end{quote}

\subsection*{You'll want to know: \zh{病人}}
\textbf{\zh{病人}} — \emph{bìngrén} — \textquotedblleft{}patient; sick person\textquotedblright{}.

\begin{quote}
\textbf{\zh{病人想喝水}} \emph{bìngrén xiǎng hē shuǐ} — The patient would like some water.
\end{quote}

The new character, in pieces: \zh{病} — outside, the sickbed piece of five strokes: a dot, a bar, a long fall down the left, then a dot and a rising stroke on the left side; inside it, five strokes: a bar, a box with a hooked right side, and a small \zh{人} in the box. The next lesson puts it in your hand.

\begin{cousinweb}
\textbf{\zh{病人}} (\emph{bìngrén}) is built from \zh{病} 'illness' and \zh{人} 'person', both inherited from Middle Chinese, \emph{bjaengH} and \emph{nyin} (Guangyun, in Baxter's transcription); the dictionary glosses the first 'grief' and 'suffering', and quotes an older book: an illness grown worse.
\end{cousinweb}

\begin{grammarlens}[title={What you've built}]
One more word for giving reasons.
\end{grammarlens}

\subsection*{Your turn}
\begin{itemize}
\raggedright
  \item \emph{Say it:} \emph{bìngrén}
  \item \emph{Say it:} \emph{bìngrén}, once more
  \item \emph{Say it:} \emph{bìngrén xiǎng hē shuǐ} — The patient would like some water.
  \item \emph{Recall:} say the Mandarin for convenient, then the Mandarin for interesting, then say \emph{bìngrén} again
\end{itemize}

\begin{tcolorbox}[breakable,colback=teal!4,colframe=teal!35!black,title={Before you move on}]
What does \zh{病人} mean? (\textquotedblleft{}Patient; sick person\textquotedblright{}.) Say it once more.
\end{tcolorbox}
\section[\texorpdfstring{\zh{病} (bìng)}{病 (bìng)}]{\zh{病} — the sickbed piece around a box}

\label{lesson:ZH-W130-bing-ill}



\begin{quote}
The word first, then the character inside it.

\begin{itemize}
\raggedright
  \item \emph{Recall:} say \emph{bìngrén} — \textbf{R1}, one lesson back, before it decays
\end{itemize}

\begin{quote}
\zh{人} \zh{男}
\end{quote}

\emph{Read it:} each one aloud, with nothing said first — \emph{rén}, \emph{nán}
\end{quote}

\subsection*{Script — the sickbed piece around a box}
\hlblockfigure{\detokenize{figures/ZH-W130-bing-ill-filmstrip.pdf}}{How \zh{病} is written, stroke by stroke}

\begin{quote}
\zh{病}
\end{quote}

\textbf{Two pieces: the sickbed piece outside, a boxed piece inside.}

\textbf{The sickbed piece, five strokes.} A dot at the top, a horizontal under it, a long left-falling stroke down the left side, then on the left a short dot and a short stroke rising to the right.

\textbf{The inside piece, five strokes.} A horizontal across the top; a short vertical, the left side of the box; a horizontal that turns down the right side and hooks; then a small \zh{人} from the top bar into the box: its left-falling stroke, then its dot.

5 + 5 = ten.

\subsection*{Your turn}
\begin{itemize}
\raggedright
  \item \emph{Write it:} \zh{病}, all ten, with the model in front of you
  \item \emph{Cover:} write it again from memory, and count the pieces, not the strokes
\end{itemize}

\begin{tcolorbox}[breakable,colback=teal!4,colframe=teal!35!black,title={Before you move on}]
How many pieces in \zh{病}? (\textbf{Two} — the sickbed piece and the boxed piece.) What sits on the left of the sickbed piece? (\textbf{A dot and a rising stroke}.) What is inside the box? (\textbf{A small \zh{人}}.)

Source: \href{https://raw.githubusercontent.com/chanind/hanzi-writer-data/68d10a4b21150cae5e1ebbd223eed289cf32d90c/data/\%E7\%97\%85.json}{Hanzi Writer Data \zh{病}, pinned stroke paths and medians}.
\end{tcolorbox}
\section[gǎnmào]{gǎnmào — to catch a cold; a cold}

\label{lesson:ZH-C130-ganmao}



\begin{quote}
Before the new one: say the Mandarin for interesting, then the Mandarin for a sick person.
\end{quote}

\subsection*{You'll want to know: gǎnmào}
\emph{gǎnmào} — \textquotedblleft{}to catch a cold; a cold\textquotedblright{}.

\begin{quote}
\emph{wǒ gǎnmào le} — I've caught a cold.
\end{quote}

Its characters come later; for now it is a word for the ear.

\begin{cousinweb}
\emph{gǎnmào} is built from \emph{gǎn} and \emph{mào}, both inherited from Middle Chinese, \emph{komX} and \emph{mawH} (Guangyun, in Baxter's transcription). The dictionary glosses the first 'to move' and the second 'to cover' and 'to wade'; how the pair came to mean a cold is not settled.
\end{cousinweb}

\begin{grammarlens}[title={What you've built}]
One more word for giving reasons.
\end{grammarlens}

\subsection*{Your turn}
\begin{itemize}
\raggedright
  \item \emph{Say it:} \emph{gǎnmào}
  \item \emph{Say it:} \emph{gǎnmào}, once more
  \item \emph{Say it:} \emph{wǒ gǎnmào le} — I've caught a cold.
  \item \emph{Recall:} say the Mandarin for interesting, then the Mandarin for a sick person, then say \emph{gǎnmào} again
\end{itemize}

\begin{tcolorbox}[breakable,colback=teal!4,colframe=teal!35!black,title={Before you move on}]
What does gǎnmào mean? (\textquotedblleft{}To catch a cold; a cold\textquotedblright{}.) Say it once more.
\end{tcolorbox}
\section[fāshāo]{fāshāo — to have a fever}

\label{lesson:ZH-C130-fashao}



\begin{quote}
Before the new one: say the Mandarin for a sick person, then the Mandarin for to catch a cold.
\end{quote}

\subsection*{You'll want to know: fāshāo}
\emph{fāshāo} — \textquotedblleft{}to have a fever\textquotedblright{}.

\begin{quote}
\emph{wǒ fāshāo le wǒ xiǎng shuìjiào} — I've got a temperature and I want to sleep.
\end{quote}

Its characters come later; for now it is a word for the ear.

\begin{cousinweb}
\emph{fāshāo} is built from \emph{fā} 'to start' and \emph{shāo} 'to burn', both inherited from Middle Chinese, \emph{pjot} and \emph{syew} (Guangyun, in Baxter's transcription).
\end{cousinweb}

\begin{grammarlens}[title={What you've built}]
One more word for giving reasons.
\end{grammarlens}

\subsection*{Your turn}
\begin{itemize}
\raggedright
  \item \emph{Say it:} \emph{fāshāo}
  \item \emph{Say it:} \emph{fāshāo}, once more
  \item \emph{Say it:} \emph{wǒ fāshāo le wǒ xiǎng shuìjiào} — I've got a temperature and I want to sleep.
  \item \emph{Recall:} say the Mandarin for a sick person, then the Mandarin for to catch a cold, then say \emph{fāshāo} again
\end{itemize}

\begin{tcolorbox}[breakable,colback=teal!4,colframe=teal!35!black,title={Before you move on}]
What does fāshāo mean? (\textquotedblleft{}To have a fever\textquotedblright{}.) Say it once more.
\end{tcolorbox}
\section[téng]{téng — sore; painful}

\label{lesson:ZH-C130-teng}



\begin{quote}
Before the new one: say the Mandarin for to catch a cold, then the Mandarin for to have a fever.
\end{quote}

\subsection*{You'll want to know: téng}
\emph{téng} — \textquotedblleft{}sore; painful\textquotedblright{}.

\begin{quote}
\emph{wǒ de tóu hěn téng} — My head really hurts.
\end{quote}

Its characters come later; for now it is a word for the ear.

\begin{cousinweb}
\emph{téng} is inherited from Middle Chinese \emph{downg} (Guangyun, in Baxter's transcription), glossed 'pain'. The regular sound change would give \emph{tóng}, not \emph{téng}; why the vowel differs is not settled.
\end{cousinweb}

\begin{grammarlens}[title={What you've built}]
One more word for giving reasons.
\end{grammarlens}

\subsection*{Your turn}
\begin{itemize}
\raggedright
  \item \emph{Say it:} \emph{téng}
  \item \emph{Say it:} \emph{téng}, once more
  \item \emph{Say it:} \emph{wǒ de tóu hěn téng} — My head really hurts.
  \item \emph{Recall:} say the Mandarin for to catch a cold, then the Mandarin for to have a fever, then say \emph{téng} again
\end{itemize}

\begin{tcolorbox}[breakable,colback=teal!4,colframe=teal!35!black,title={Before you move on}]
What does téng mean? (\textquotedblleft{}Sore; painful\textquotedblright{}.) Say it once more.
\end{tcolorbox}
\section[\texorpdfstring{\zh{看病} (kànbìng)}{看病 (kànbìng)}]{\zh{看病} (kànbìng) — to see a doctor}

\label{lesson:ZH-C130-kanbing}



\begin{quote}
Before the new one: say the Mandarin for to have a fever, then the Mandarin for sore.
\end{quote}

\subsection*{You'll want to know: \zh{看病}}
\textbf{\zh{看病}} — \emph{kànbìng} — \textquotedblleft{}to see a doctor\textquotedblright{}.

\begin{quote}
\textbf{\zh{我病了我去看病}} \emph{wǒ bìng le wǒ qù kànbìng} — I'm ill, so I'm off to the doctor.
\end{quote}

\begin{cousinweb}
\textbf{\zh{看病}} (\emph{kànbìng}) is built from \zh{看} 'to look at' and \zh{病} 'illness', both inherited from Middle Chinese, \emph{khanH} and \emph{bjaengH} (Guangyun, in Baxter's transcription): to have an illness looked at.
\end{cousinweb}

\begin{grammarlens}[title={What you've built}]
One more word for giving reasons.
\end{grammarlens}

\subsection*{Your turn}
\begin{itemize}
\raggedright
  \item \emph{Say it:} \emph{kànbìng}
  \item \emph{Say it:} \emph{kànbìng}, once more
  \item \emph{Say it:} \emph{wǒ bìng le wǒ qù kànbìng} — I'm ill, so I'm off to the doctor.
  \item \emph{Recall:} say the Mandarin for to have a fever, then the Mandarin for sore, then say \emph{kànbìng} again
\end{itemize}

\begin{tcolorbox}[breakable,colback=teal!4,colframe=teal!35!black,title={Before you move on}]
What does \zh{看病} mean? (\textquotedblleft{}To see a doctor\textquotedblright{}.) Say it once more.
\end{tcolorbox}
